\section{Shift comparison}
\label{shift:section}

Shift comparison lowers the degree of a positive test after multiplying
by a translate of an arbitrary nonnegative function.  The degree-one
case uses Bohr-set almost-periodicity.  At higher degree we recover the
test from two fixed finite lists and apply the induction outside one
exceptional set of shifts.  The lower-degree test may then depend on
the shift.

Put $G=\Z/N\Z$, with its uniform probability measure, and evaluate
nilsequences on the representatives $0,\ldots,N-1$.
For $d\geq 0$ and $q\geq2$, let $\mathcal T_d(q)$ denote the
$[0,1]$-valued niltests of degree at most $d$ and complexity at most $q$.
Thus $\mathcal T_0(q)$ consists of constants in $[0,1]$.

\begin{theorem}[Shift comparison]
\label{shift:required}
For fixed $d\geq1$ and $0<\epsilon<1$, there is a constant
$C=C(d,\epsilon)$ such that the following holds.  Suppose that
$p\geq2$, $N\geq\exp((2+p)^C)$ is odd, and
$0\leq f,g,J\leq e^p$.  Set $P=(2+p)^C$.  If
\begin{equation}\label{shift:upper-hypothesis}
 \E_n(f-g)(n)T(n)\leq e^{-P}
 \qquad(T\in\mathcal T_d(P)),
\end{equation}
then there is a set $E\subseteq G$, depending on $f,g,J$, with
$|E|\leq e^{-p}N$, such that
\begin{equation}\label{shift:required-conclusion}
 \E_n\bigl(f-(1+\epsilon)g\bigr)(n)J(n+h)T'(n)\leq e^{-p}
 \qquad(h\notin E,\ T'\in\mathcal T_{d-1}(p)).
\end{equation}
The exceptional set is chosen before $T'$.
\end{theorem}

The proof is given in Proposition~\ref{shift:degree-one} for $d=1$
and Proposition~\ref{shift:higher} for higher degrees.  The latter
uses the structural results of Section~\ref{prelim:section} and
Lemma~\ref{prelim:step-drop}.

\subsection{The degree-one argument}

For a finite set $\Omega$ of characters of $G$ and $0<w\le1$, put
\[
 B(\Omega,w)=\{x\in G:|\chi(x)-1|\le w\text{ for every }\chi\in\Omega\}.
\]
We use these symmetric common-width Bohr sets, retaining $\Omega$ and
$w(B)=w$ as part of the data; their rank is $D=|\Omega|$.
The whole group has the presentation with one trivial character and
width one.  A regular Bohr set $B$ of rank $D\geq1$
satisfies
\[
 |B_{1+t}|=(1+O(D|t|))|B|\qquad (|t|\leq c/D).
\]
Here dilating a Bohr set means dilating its defining widths.
The usual covering argument gives
$|B|/N\geq(w(B)/10)^D$ and
$|B_\rho|\geq(\rho/C)^{CD}|B|$ for $0<\rho<1$; a regular
dilate exists between any two scales whose ratio is two.
Also, if $v\in B_\rho$, regularity gives
\begin{equation}\label{shift:bohr-tv}
 \|\mu_B(\cdot-v)-\mu_B\|_1\ll D\rho,
\end{equation}
Here $\mu_B(x)=1_B(x)/|B|$ is the probability mass function on $B$,
and the norm in~\eqref{shift:bohr-tv} is the unnormalized sum over $G$.
Regular dilation and the following almost-periodicity input preserve
the common-width convention.  Widths will be tracked explicitly;
adding characters does not permit a width-independent estimate for
the relative size of the resulting Bohr set.

The convolution almost-periodicity method originates with Croot and
Sisask~\cite{CrootSisask2010}. We use the following radius-sensitive
Bohr-set form of Schoen and Sisask.

\begin{lemma}[Almost-periodicity with a width bound]\label{shift:ap-width}
Fix $0<\eta<1/2$.  If $S,S'$ are regular of rank $D\geq1$,
$A_1\subseteq S$, $A_2\subseteq s_0+S'$, both relative densities are
at least $e^{-q}$ with $q\geq1$, and $|K|\leq2|S|$, then, for independent
uniform $a_i\in A_i$, the probability that $a_1-a_2$ belongs to $K$
changes by at most $\eta$ after adding an independent uniform element
of a regular Bohr set $D_*\subseteq S'$ with
\begin{equation}\label{shift:ap-cost}
 \rank D_*\leq D+C_\eta(1+q)^4,\qquad
 w(D_*)\geq w(S')\exp\bigl(-C_\eta(1+q+\log(2+D))\bigr).
\end{equation}
\end{lemma}
\begin{proof}
We apply Schoen--Sisask \cite[Theorem~5.4]{SchoenSisask2016}, which
uses unnormalized counting convolution.  In its notation, a regular
Bohr set of rank $d$ and radius $\rho$, a subset of relative density
greater than $\sigma$, a sumset bound $|A+S_0|\leq K_1|A|$, and
$\beta=|M|/|L|\leq1$ give a regular sub-Bohr set of rank at most
$d+d'$ and radius at least $c\rho\eta\sqrt\beta/(d^2d')$, where
\[
 d'\leq C\eta^{-2}\log^2(2/(\eta\beta))
       \log(2/\beta)\log(2K_1)+C\log(1/\sigma).
\]
The error in its three-fold convolution is at most $\eta|A||M|$.

Write $\alpha_i$ for the two relative densities, and take a regular
$B_0=S'_\tau$ with $\tau\asymp1/D$ sufficiently small.  Use $B_0$
as the ambient Bohr set, and set $A=-A_2$, $S_0=B_0$.
The increment set $S_0$ has density one in $B_0$, so
$\sigma=1/2$ is valid; regularity gives
$|A+S_0|\leq2|S'|\leq(2/\alpha_2)|A|$.
If $|K|\geq|A_1|$, take $M=A_1$, $L=-K$; then
$\alpha_1/2\leq\beta\leq1$.
If $\eta|A_1|\leq|K|<|A_1|$, interchange $M$ and $L$;
then $\eta\leq\beta\leq1$, and the error is still at most
$\eta|A_1||A_2|$.
In either case $d'\leq C_\eta(1+q)^4$, and the radius lower bound
implies \eqref{shift:ap-cost}.  The counting convolution at zero,
divided by $|A_1||A_2|$, is the probability in the statement;
symmetry of $D_*$ gives the asserted smoothing estimate.
Finally, if $|K|<\eta|A_1|$, every translate of the difference law
assigns $K$ probability at most $|K|/|A_1|<\eta$.
Taking $D_*=B_0$ proves this case as well.
\end{proof}

Only fixed $\eta>0$ is used here.  All convolution normalizations in
the following local calculation are stated explicitly.

\begin{lemma}[A bilinear reduction]\label{shift:bilinear}
Let $a=f-(1+\epsilon)g$.  In degree one it suffices to prove
\begin{equation}\label{shift:bilinear-bound}
 \E_{x,y}H(x)W(y)a(x+y)\leq e^{-4p}
 \qquad(0\leq H,W\leq1).
\end{equation}
\end{lemma}
\begin{proof}
If the conclusion fails on a set $E$ of more than $e^{-p}N$ shifts,
then $\E_n a(n)J(n+h)>e^{-p}$ for every $h\in E$: a constant test
cannot create a positive violation where this inequality fails.
Take $H= e^{-p}J$ and $W(y)=1_E(-y)$.  The left side of
\eqref{shift:bilinear-bound} equals
$e^{-p}\E_h1_E(h)\E_n a(n)J(n+h)>e^{-3p}$, a contradiction.
\end{proof}

\paragraph{The refinement to be proved.}
To prove the bilinear estimate, we decompose its two independent
averages into smaller Bohr cells. The difficulty is that a restriction
can concentrate either weight. We measure both weights by a concave
potential, so that concentration and genuine smoothing each pay for
the further restrictions they require.

Write $a=f-(1+\epsilon)g$. A node consists of a regular Bohr set
$B$, two translates $x_B+B,y_B+B$, and functions $H,W\in[0,1]$
supported on these translates. Define
\[
 I=\mathbb E_{x\in x_B+B,\,y\in y_B+B}H(x)W(y)a(x+y),
 \enspace u=\mathbb E_{x\in x_B+B}H(x),
 \enspace v=\mathbb E_{y\in y_B+B}W(y),
 \enspace \Phi=(uv)^{1/4}.
\]
Every node average uses the full size of its cell, even after a weight
has been restricted and extended by zero. In particular,
$I\le e^puv\le e^p\Phi$. The original bilinear average is the node
with $B=G$.

We shall construct a refinement with successor nodes indexed by
$\nu$ and nonnegative coefficients $c_\nu$ such that
\begin{equation}\label{shift:one-round-refinement}
 I\le\sum_\nu c_\nu I_\nu+\varepsilon_{\rm err}\Phi,
 \qquad
 \sum_\nu c_\nu\Phi_\nu\le(1-\gamma)\Phi.
\end{equation}
Here $\gamma>0$ depends only on $\epsilon$, and
$\varepsilon_{\rm err}$ can be any sufficiently small fixed multiple
of $e^{-10p}$. Terms which need no further refinement are included
in the error in the first inequality. The same construction controls
the geometry: if $D$ and $w$ are the parent's stored rank and width,
every successor has
\begin{equation}\label{shift:one-round-geometry}
 D_\nu\le D+(2+p)^{A_\epsilon},\qquad
 \log\frac{w}{w_\nu}\le(2+p+D)^{B_\epsilon}.
\end{equation}
We need this refinement only for descendants of the root during
$R=\lceil c_{\rm depth}p\rceil$ rounds. The geometric bounds above
forecast polynomial upper bounds for the ranks and $\log(1/w)$ of
these nodes and their intermediate dilates. Choose $P$ to cover those
bounds; the refinement is then uniform over every center and pair of
weights encountered in the run.

There are two reasons to stop part of a node. If either cell mean is
tiny, its entire positive contribution is small relative to $\Phi$.
If a localized convolution is close to $1$ in a sufficiently high
moment, positive comparison with a smooth Bohr test makes that
localized contribution small. The other parts become successors.
A highly concentrated restriction is paid for by the potential lost
when that restriction is removed. In the absence of such concentration,
a large convolution moment yields almost-periods and forces a fixed
decrease of the average potential. The following two elementary
estimates quantify these last two mechanisms.

\begin{lemma}[Potential estimates]\label{shift:potential}
Put $\theta=1/4$.  The following estimates hold.
\begin{enumerate}
\item If $u,v>0$, $0\leq bu_1\leq u$, and $u_1\geq Ku$, then
\begin{equation}\label{shift:peeling}
 (uv)^\theta-((u-bu_1)v)^\theta
 \geq \theta K^{1-\theta}b(u_1v)^\theta.
\end{equation}
\item If $0\leq X\leq M$, $\E X\leq1$, and
$\E X^2\geq1+c$, where $M\geq1$, then
\begin{equation}\label{shift:concavity}
 \E\sqrt X\leq1-\frac{c}{2(1+\sqrt M)^2}.
\end{equation}
Consequently, if $X,Y\geq0$, $\E X,\E Y\leq1$, and $X$ satisfies
these additional bounds, then
$\E(XY)^{1/4}\leq1-c/[4(1+\sqrt M)^2]$.
\end{enumerate}
\end{lemma}
\begin{proof}
Concavity gives $1-(1-t)^\theta\geq\theta t$ for $0\leq t\leq1$.
Apply this with $t=bu_1/u$, and use
$(u_1/u)^{1-\theta}\geq K^{1-\theta}$.
For the second assertion,
\[
 \frac{1+x}{2}-\sqrt{x}
 =\frac{(x-1)^2}{2(1+\sqrt{x})^2}
 \geq\frac{(x-1)^2}{2(1+\sqrt M)^2}.
\]
Moreover $\E(X-1)^2\geq c$ and $(1+\E X)/2\leq1$.
The final assertion follows from Cauchy--Schwarz and
$\E\sqrt Y\leq1$.
\end{proof}

The local moment argument and the weighted dependent-random-choice
step below adapt the moment and sifting methods of Kelley and
Meka~\cite[Sections~4--5]{KelleyMeka2023}, as developed in
Bloom and Sisask~\cite[Sections~2.1--2.2]{BloomSisaskKelleyMeka}.
Here they are used inside the refinement of two weighted Bohr cells.

\begin{lemma}[Local convolution alternatives]\label{shift:local-moments}
Let $L,S$ be symmetric regular Bohr sets, with $S$ a sufficiently
small dilate of $L$.  Let $F,G\geq0$ be supported on $L$, have
$|L|^{-1}\sum F=|L|^{-1}\sum G=1$, and satisfy $F,G\leq M$.
Define
\[
 (F*G)(t)=\frac1{|L|}\sum_xF(x)G(t-x),\qquad
 (F\circ G)(t)=\frac1{|L|}\sum_xF(x+t)G(x).
\]
Write $\nu_+$ for the law of $s_1+s_2$ and $\nu_-$ for the law
of $s_1-s_2$, where $s_1,s_2$ are independent and uniform on $S$.
For fixed $\delta>0$, and integer $m\geq1$, either
\begin{equation}\label{shift:flat}
 \|F*G-1\|_{L^{2m}(\nu_+)}\leq\delta,
\end{equation}
or, for one of $D=F,G$ and an integer $m'\leq C_\delta m$,
\begin{equation}\label{shift:unbalanced}
 \|D\circ D\|_{L^{2m'}(\nu_-)}\geq1+c_\delta.
\end{equation}
It is enough that the regularity errors, multiplied by $M$, be
sufficiently small in terms of $\delta$.
\end{lemma}
\begin{proof}
For real $A,B$ and every integer $j\geq1$, expansion and
Cauchy--Schwarz give
\begin{equation}\label{shift:moment-cs}
 \|A*B\|_{L^{2j}(\nu_+)}
 \leq
 \|A\circ A\|_{L^{2j}(\nu_-)}^{1/2}
 \|B\circ B\|_{L^{2j}(\nu_-)}^{1/2}.
\end{equation}
Indeed, after expanding the $2j$th power, apply Cauchy--Schwarz
to the two factors
$\E_{s\in S}\prod_{i=1}^{2j}A(s-t_i)$ and
$\E_{s\in S}\prod_{i=1}^{2j}B(s+t_i)$, summing each $t_i$
over the group with weight $|L|^{-1}$.
The same expansion shows
\begin{equation}\label{shift:positive-moments}
 \E_{\nu_-}(A\circ A)^j\geq0\qquad(j\geq1),
\end{equation}
because the resulting expression is a sum of squares.

Put $F_1=F-1_L$.  Regularity and the cap $M$ show, uniformly on
$S+S$ and $S-S$, respectively, that
\[
 1_L*G=1+o_\delta(1),\qquad
 F\circ F=1+F_1\circ F_1+o_\delta(1).
\]
If \eqref{shift:flat} fails, then
$\|F_1*G\|_{L^{2m}(\nu_+)}\geq\delta/2$.
Either $\|G\circ G\|_{L^{2m}(\nu_-)}\geq2$, or
\eqref{shift:moment-cs} gives
$\|F_1\circ F_1\|_{L^{2m}(\nu_-)}\geq\delta^2/8$.
In the latter case put $X=F_1\circ F_1$ and choose
$m'=Rm$ with a sufficiently large integer $R=R(\delta)$.
All terms in the binomial expansion have nonnegative expectations
by \eqref{shift:positive-moments}; hence
\[
 \E(1+X)^{2m'}
 \geq \binom{2m'}{2m}\E X^{2m}
 \geq(R\delta^2/8)^{2m}.
\]
Choose $R$ so that $(R\delta^2/8)^{1/R}>1+2c_\delta$, and absorb
the uniform regularity error using the triangle inequality.
\end{proof}

\begin{proposition}[Degree-one shift comparison]\label{shift:degree-one}
Theorem~\ref{shift:required} holds for $d=1$.
\end{proposition}
\begin{proof}
We construct the refinement in \eqref{shift:one-round-refinement},
with the geometry in \eqref{shift:one-round-geometry} and the full-cell
averaging convention fixed above. Put $\theta=1/4$.

Fix $K$ so large that $\theta K^{1-\theta}>2$, and choose
$c_0>0$ with $(c_0K)^\theta\leq1/2$.  Choose
$\kappa=e^{-c_{\rm cut}p}$ and
$\zeta=\exp(-(2+p)^{C_\epsilon'})$ sufficiently small.
Choose $c_{\rm cut}=c_{\rm cut}(\epsilon)$ first to control the
tiny-node error below.  The moment constant and the number of rounds
will be chosen afterwards; all are independent of the node rank.
If $\min(u,v)<\kappa$, stop the node, since its positive
contribution is at most
\begin{equation}\label{shift:tiny-node}
 e^puv\leq e^p\kappa^{1-\theta}\Phi.
\end{equation}

Here and throughout this proof, the required scale choices are regular
dilates whose logarithmic width losses are polynomial in $p+D$, where
$D=\rank B$.  Successor ranks are at most
$D+(2+p)^{C_\epsilon''}$.  To make the set of admissible successor
shapes precise, retain $w=w(B)$ and allow all regular presentations $C$
satisfying
\[
 C\subseteq B_{c\zeta/(2+D)},\qquad
 \rank C\leq D+(2+p)^{A_\epsilon},\qquad
 w(C)\geq w\exp(-(2+p+D)^{B_\epsilon}).
\]
Choose $A_\epsilon$ from \eqref{shift:ap-cost}, which will be
applied with $q=O_\epsilon(p^2)$.  Choose $B_\epsilon$ after
the exponent defining $\zeta$, large enough to cover that
almost-periodicity width loss and the bounded number of regular
dilations within one refinement.  These choices
precede peeling.  The exponent defining $P$ is chosen last.

\emph{Removing concentration.}
While a translate of an admissible $C$ has mean $u_1>Ku$ of $H$,
remove the whole slice from $H$.  Set $b=|C|/|B|$.
The removed contribution is represented, up to
$O(e^p\zeta b)$, by nodes of shape $C$ with coefficient $b$:
keep the first center fixed, average the second center in its parent
cell, and add an independent uniform element of $C$ on that side.
This error bound is \eqref{shift:bohr-tv}, applied to the second
parent cell.  The second means average to at most $v$.
Consequently their total potential is at most $b(u_1v)^\theta$.
Inequality~\eqref{shift:peeling} pays a factor greater than two for
this potential.  Perform the same operation if $W$ has such a slice.

Stop immediately if either current mean becomes smaller than $\kappa$.
At every other removal, the removed mass is at least
$K\kappa b_{\min}>0$, where the volume bound and admissibility give
\[
 b\geq |C|/N\geq b_{\min}:=
 \left(\frac{w\exp(-(2+p+D)^{B_\epsilon})}{10}\right)^{D+(2+p)^{A_\epsilon}}.
\]
Thus peeling terminates in finitely many steps.  Moreover its errors
are summable without multiplying by the number of steps: since
$u_1,v\geq\kappa$, each $b$ is at most
$\kappa^{-2\theta}b(u_1v)^\theta$, and
\eqref{shift:peeling} telescopes.
It remains to treat a non-tiny remainder for which every admissible
translate has $H$-mean at most $Ku$ and $W$-mean at most $Kv$.

\emph{Matching two local scales.}
Choose tiny regular dilates $S\ll L\ll B$. Choose $x_0,y_0$
independently and uniformly in the two parent cells, take
$r\in L$ and $s=s_1+s_2$ with $s_1,s_2\in S$ independently
uniform, and put
\[
 x=x_0+r,\qquad y=y_0-r+s.
\]
The resulting average differs from the parent average by
$O(e^p\zeta)$, by \eqref{shift:bohr-tv} twice.  Restrict both
weights to $x_0+L$ and $y_0+L$; the further error is
$O(e^p\zeta)$ because $s$ is on a much smaller scale than $L$.
Let $u_0,v_0$ be the resulting means on $L$.
The absence of concentration gives $u_0\leq Ku$, $v_0\leq Kv$;
independence of $x_0,y_0$ and concavity give
\begin{equation}\label{shift:matched-potential}
 \E_{x_0,y_0}(u_0v_0)^\theta\leq(uv)^\theta.
\end{equation}

One can return to independent cells of any sufficiently small common
shape $C$ by adding independent $t_1,t_2\in C$ to $x,y$.
In $(r,s)$ coordinates this translates the law by
$(t_1,t_1+t_2)$, so its total variation changes by $O(\zeta)$.
If $U,V$ are the new cell means, then, exactly,
\begin{equation}\label{shift:means-exact}
 \E_{r,s}U\leq u_0,\quad \E_{r,s}V\leq v_0,
 \quad \E_{r,s}(UV)^\theta\leq(u_0v_0)^\theta.
\end{equation}
For example, for each fixed $s,t_2$, summing the zero-extended second
weight over $-L+s+t_2$ is at most its sum over the whole group,
namely $|L|v_0$.  This proves the second inequality; the first is
identical and the last follows by Cauchy--Schwarz and concavity.

If $u_0<c_0u$ or $v_0<c_0v$, return to independent cells in this
way.  The total potential of these centers is at most
$(c_0K)^\theta\Phi\leq\Phi/2$.
For the other centers $u_0,v_0\geq c_0\kappa$.  Fix one and
normalize its truncated weights to functions $F,G$ on $L$ of mean
one.  Their caps are at most $M=(c_0\kappa)^{-1}$.
Choose $m=\lceil c_{\rm mom}p\rceil$ with
$c_{\rm mom}=c_{\rm mom}(\epsilon,c_{\rm cut})$ large enough, and apply
Lemma~\ref{shift:local-moments} with sufficiently small fixed
$\delta=\delta(\epsilon)$.

\emph{The flat alternative terminates a matched term.}
Its integral, divided by $u_0v_0$, is
$\E_{s\sim\nu_+}a(x_0+y_0+s)(F*G)(s)$.
Outside a set of $\nu_+$-measure at most $2^{-2m}$,
$|F*G-1|\leq2\delta$.  Since $0\leq F*G\leq M$, the exceptional
part costs at most $O(e^p(M+1)2^{-2m})$.
On the complementary set use
\[
 a(F*G)\leq (1+2\delta)f
              -(1+\epsilon)(1-2\delta)g.
\]
Choose $\delta$ so that
$(1+\epsilon)(1-2\delta)\geq1+2\delta$.

Hypothesis~\eqref{shift:upper-hypothesis} gives
$\E_{s\sim\nu_+}(f-g)(x_0+y_0+s)\leq\eta$
for any prescribed $\eta=\exp(-(2+p)^{O_\epsilon(1)})$,
after taking $P$ sufficiently large.  Here is the needed
approximation justification.  Approximate a translated regular Bohr
indicator from above and below by $[0,1]$-valued Lipschitz functions
on its defining torus, with transition width $\rho w(S)$.  Their
difference has normalized group mean at most $O(D\rho)|S|/N$, and
their Lipschitz constants are at most $O(D/(\rho w(S)))$.
These are degree-one niltests: the torus has lattice $\Z^D$ and
the orbit is given by the defining characters.  With
$\alpha=|S|/N$, the hypothesis and the signed comparison yield
\[
 \E_{s\in S}(f-g)(z+s)
 \leq \alpha^{-1}e^{-P}+O(e^pD\rho)
 \qquad(z\in G).
\]
The width and rank bookkeeping below bounds both $\log(1/w(S))$
and $\log(1/\alpha)$ polynomially in $p$.  Thus $\rho$ can be chosen
with polynomial logarithmic loss to make the last display at most
$\eta$, with all test complexities at most $P$.
Averaging over the second copy of $S$ proves the assertion for $\nu_+$.
Thus this node has positive contribution at most an arbitrarily small
multiple of $e^{-10p}(u_0v_0)^\theta$.

\emph{The unbalanced alternative contracts the potential.}
Let $Q$ be the function supplied by \eqref{shift:unbalanced}, and
write $q=2m'=O_\epsilon(p)$.
Choose a tiny regular shrink $S'$ of $S$.  Regularity, with error
small after multiplication by $M^{2q}$, supplies an $s_0\in S$ such that
\[
 \E_{a\in S,\ b\in s_0+S'}(Q\circ Q)(a-b)^q
 \geq(1+c)^q
\]
for a fixed $c>0$.  Set
$K_0=\{Q\circ Q\geq1+c/2\}\cap(S-s_0-S')$.
Regularity gives $|K_0|\leq2|S|$.

For dependent random choice with these real weights, let $L^+$ be a
slight dilate of $L$ containing $L-a$ for every
$a\in S\cup(s_0+S')$.  The scale separation and regularity ensure
$|L^+|\leq2|L|$.  Independently sample $t_1,\ldots,t_q$
uniformly from $L^+$.  Given these samples, include each
$a\in S$ in $A_1$, and each $b\in s_0+S'$ in $A_2$,
independently, with probabilities
\[
 \prod_{j=1}^q\frac{Q(t_j+a)}{M},\qquad
 \prod_{j=1}^q\frac{Q(t_j+b)}{M}.
\]
The expected pair weight is
$\bigl(|L|/(|L^+|M^2)\bigr)^q(Q\circ Q)(a-b)^q$.
The contribution of pairs outside $K_0$, relative to the lower
bound for all pairs, is at most
$((1+c/2)/(1+c))^q$.  By increasing $m$ by a fixed factor,
this is smaller than any chosen fixed constant.
To justify the simultaneous density and bad-pair bounds, write
$Z=|A_1||A_2|/(|S||S'|)$ and let $Z_{\rm bad}$ be the same
normalized count restricted to differences outside $K_0$.
The preceding estimates give $\E Z\ge\exp(-O_\epsilon(p^2))$
and $\E Z_{\rm bad}\le\delta_*\E Z$, where $\delta_*$ can
be made any sufficiently small fixed constant by increasing the
fixed multiple defining $q$.  Given a sufficiently large fixed
$C_*$, take $C_*\delta_*\le1/2$.  Some outcome then has
\[
 Z-C_*Z_{\rm bad}\ge\tfrac12\E Z.
\]
It follows that $Z\ge\exp(-O_\epsilon(p^2))$ and
$Z_{\rm bad}/Z\le C_*^{-1}$.  Each relative density is at least
$Z$, since both are at most one.  Thus both densities have the
required lower bound and an arbitrarily high fixed fraction of
the differences lie in $K_0$.

Apply \eqref{shift:ap-cost} to these sets.  Choose a tiny regular
shrink $C$ of $D_*$ so that adding independent opposite shifts
from $C$ changes the uniform law on $D_*$ by $O(\zeta)$.
If $d$ is uniform on $D_*$ and $c_1,c_2$ are independent and
uniform on $C$, the almost-periodicity estimate therefore gives
\[
 \E_{a\in A_1,b\in A_2,d,c_1,c_2}
 1_{K_0}(a-b+d+c_1-c_2)\geq1-\eta_0,
\]
where $\eta_0>0$ is any sufficiently small fixed constant.
For a finite set $A$, write
$\mathsf P_Av(x)=|A|^{-1}\sum_{t\in A}v(x+t)$; these averaging
operators do not use the $|L|^{-1}$ normalization of function-function
convolution.  Since $Q\circ Q\geq1+c/2$ on $K_0$ and is
nonnegative everywhere, the last display implies
\[
 \max_x\mathsf P_C\mathsf P_{-C}(Q\circ Q)(x)
 \geq(1+c/2)(1-\eta_0)\geq1+c'
\]
for a fixed $c'>0$.
This function is the autocorrelation of $\mathsf P_CQ$, with sums
normalized by $|L|$.  Cauchy--Schwarz puts its maximum at zero.
Its mass outside $L$ contributes $O(M^2\zeta)$ to the second
moment, so
$\E_{r\in L}(\mathsf P_CQ)(r)^2\geq1+c'/2$.
The same estimate holds on the reflected and slightly translated
law used by the second weight, by regularity.

Return to independent cells of shape $C$, with centers $x_0+r$
and $y_0-r+s$, where $r$ is uniform on $L$ and $s\sim\nu_+$
independently.  Their means satisfy
\[
 U/u_0=\mathsf P_CF(r),\qquad
 V/v_0=\mathsf P_CG(-r+s).
\]
This shape is admissible: its rank increase is polynomial in $p$,
and its logarithmic width loss relative to $B$ is polynomial in
$p+D$, by \eqref{shift:ap-cost} and the intervening regular dilations.
The absence of concentration gives
$U/u_0,V/v_0\leq K/c_0$.  Equations~\eqref{shift:means-exact}
give means at most one, and one of these variables has second
moment at least $1+c'/4$.
Lemma~\ref{shift:potential} therefore contracts the local potential
by a fixed factor $q_0<1$.

\emph{Summing the refinements.}
Let $A$ be the matched potential from small-ratio centers and $B_1$
that from the other centers.  We have $A\leq\Phi/2$ and
$A+B_1\leq\Phi$, so the surviving potential is at most
$A+q_0B_1\leq(1+q_0)\Phi/2$.
For completeness, let $S_{\rm peel}$ be the total weighted
potential already extracted by peeling, and $R_{\rm rem}$ the
potential of its remainder.  Telescoping gives
$2S_{\rm peel}+R_{\rm rem}\le\Phi$.  The matching argument
contracts the surviving remainder potential by
$\lambda=(1+q_0)/2<1$; a tiny remainder has no successors.
Hence the total successor potential is at most
\[
 S_{\rm peel}+\lambda R_{\rm rem}
 \le\max(1/2,\lambda)\Phi.
\]
This gives a fixed $\gamma=\gamma(\epsilon)>0$ such that all
successors of a node have total weighted potential at most
$(1-\gamma)\Phi$.
Their integrals bound its integral from above up to
$c_\epsilon e^{-10p}\Phi$, where $c_\epsilon$ can be made
arbitrarily small.  Indeed all errors before a tiny-mean cutoff
are bounded by fixed powers of $e^p\kappa^{-1}$ times $\zeta$;
the extracted errors have already been summed by
\eqref{shift:peeling}.

This proves \eqref{shift:one-round-refinement} with
$\varepsilon_{\rm err}=c_\epsilon e^{-10p}$; admissibility of the
successor shapes gives \eqref{shift:one-round-geometry}.

Apply the refinement for $R=\lceil c_{\rm depth}p\rceil$ rounds.
After $r$ rounds, the total weighted potential of the remaining nodes
is at most $(1-\gamma)^r\Phi_{\rm root}$. Hence the sum of all
refinement errors is at most
$\varepsilon_{\rm err}\Phi_{\rm root}/\gamma$, and the remaining
positive integral is at most
\[
                    e^p(1-\gamma)^R\Phi_{\rm root}.
\]
Since $\Phi_{\rm root}\le1$, choosing $c_{\rm depth}$ sufficiently
large and then the permitted error constant sufficiently small makes
their sum less than $e^{-4p}$.

It remains to justify one comparison budget for the whole refinement.
Along every branch, \eqref{shift:one-round-geometry} gives
\[
 D_{r+1}\le D_r+(2+p)^{A_\epsilon},\qquad
 \log(1/w_{r+1})\le\log(1/w_r)+(2+p+D_r)^{B_\epsilon}.
\]
Starting with $D_0=1,w_0=1$, both $D_r$ and $\log(1/w_r)$ are
polynomial in $p$ for $r\le R$. The same is true of the intermediate
dilates used inside a round. The Bohr volume bound
$|B|/N\ge(w(B)/10)^D$ consequently gives a polynomial bound for
$\log(N/|B|)$. One choice $P=(2+p)^C$ therefore covers the
smooth-test approximations, their required errors, and the size
cutoff at every node. This proves the bilinear estimate and, by
Lemma~\ref{shift:bilinear}, degree-one shift comparison.
\end{proof}

\subsection{The finite-list reduction}

For the induction from degree $r$ to degree $r+1$, a violating
positive test of degree $r$ may be chosen separately at every shift.
We cannot take a union bound over all these tests. Instead we will
recover them, on a sufficiently large set of shifts, from two fixed
finite lists of positive tests: one observed at $n$, the other at
$n+h$. The remaining factor may depend on $h$, but has degree at
most $r-1$. Degree-$r$ shift comparison is uniform over this remaining
factor, so only the two fixed lists enter the union bound.
The next lemma gives the precise reduction. The rest of the proof
constructs its factorization.

\begin{lemma}[Why fixed lists suffice]\label{shift:fixed-lists-suffice}
Assume shift comparison at degree $r\geq1$.
Let $0<\epsilon<1$, $0\leq f,g,J\leq e^p$, and suppose that
$N$ is odd and $N\geq\exp(P)$ for the sufficiently large
polynomial budget $P$ in the conclusion.  Suppose, for a set
$H\subseteq G$ of density at least $\sigma$, that positive tests
$T_h$ satisfy
\[
 \E_n\bigl(f-(1+\epsilon)g\bigr)(n)J(n+h)T_h(n)>\delta.
\]
Suppose there are fixed families $A_i,B_j\in\mathcal T_d(Q)$,
$d=r+1$, of sizes at most $F$, and, for each $h\in H$,
an approximation
\[
 T_h(n)=\sum_{\ell=1}^{M}A_{i(\ell,h)}(n)
 B_{j(\ell,h)}(n+h)C_{\ell,h}(n)+e_h(n),
 \qquad C_{\ell,h}\in\mathcal T_{r-1}(Q),
\]
with $\E|e_h|\leq\delta e^{-2p}/6$.
Here all factors take values in $[0,1]$.
If
\[
 Q+\log(F+M+2)+\log(1/\sigma)+\log(1/\delta)
 \leq(2+p)^{O_{d,\epsilon}(1)},
\]
then \eqref{shift:upper-hypothesis}, with $P$ a sufficiently
large polynomial in $p$, rules out these violations.
\end{lemma}
\begin{proof}
Choose $q\geq\max(p,Q)$ polynomially large so that
$F^2e^{-q}<\sigma$ and $Me^{-q}<\delta/2$.
For each fixed $i,j$, apply degree-$r$ shift comparison to
$fA_i,gA_i,JB_j$, with multiplicative slack $\epsilon/2$ and
budget $q$.  Its hypothesis follows from
\eqref{shift:upper-hypothesis}: multiplying a degree-$r$ test
by $A_i$ gives a degree-$d$ positive test, of complexity bounded
polynomially in the required budgets.
There are at most $F^2$ fixed pairs, so some $h\in H$ lies
outside every exceptional set.  Uniformity over all lower tests
then bounds each summand by $e^{-q}$.  Replacing
$f-(1+\epsilon/2)g$ by $f-(1+\epsilon)g$ only decreases the
integral.  Finally
$|(f-(1+\epsilon)g)J|\leq3e^{2p}$, so the approximation error
costs at most $\delta/2$.  This contradicts the displayed
violation.
\end{proof}

\subsection{Modeling and a common anchor}

To construct the fixed lists, we first approximate each input by a
finite linear combination of niltests in a seminorm adapted to the
shifted comparisons. The separation argument is a form of the
Hahn--Banach method for deriving decompositions from inverse theorems
\cite[Section~3.2]{GowersHahnBanach2010}; the explicit finite list is
then obtained by sampling.  On a finite probability space, the extended dual of a
seminorm $X$ is
\[
 X^*(\psi)=\sup_{X(v)\le1}|\langle v,\psi\rangle|\in[0,\infty].
\]
It is finite precisely when $\psi$ annihilates $\ker X$.
A family is \emph{balanced} if it is closed under multiplication by
complex scalars of modulus at most one.

\begin{lemma}[Modeling from a seminorm detector]\label{shift:modeling}
Let $\Omega$ be a finite probability space, let $X$ be a seminorm
on $\C^\Omega$ with $X(v)\leq K\|v\|_1$, and let
$\mathcal S$ be a nonempty compact, balanced collection of functions of
supremum at most one.  Fix $M,K\geq1$, $\beta>0$, and $0<\tau<1$.
Suppose that every $\psi$ satisfying
\[
 \|\psi\|_\infty\leq K/\tau,
 \qquad X(\psi)\geq\tau/M^2
\]
has $|\langle\psi,Q\rangle|\geq\beta$ for some $Q\in\mathcal S$.
Then every $b$ with $\|b\|_\infty\leq M$ has a representation
\[
 b=\sum_{i=1}^{\ell}c_iQ_i+e,\qquad
 \sum_i|c_i|\leq2/\beta,
 \qquad X(e)\leq2\tau,
\]
with $\ell\leq1+4K^2/(\beta^2\tau^2)$.
\end{lemma}
\begin{proof}
Write $R=2/\beta$.  Suppose that $b$ does not belong to
$R\operatorname{conv}(\mathcal S)+\{e:X(e)\leq\tau\}$.
Separation in the underlying finite-dimensional real vector space
gives a nonzero $\psi$ with
\[
 \operatorname{Re}\langle b,\psi\rangle>
 R\sup_{Q\in\mathcal S}|\langle Q,\psi\rangle|
 +\tau X^*(\psi).
\]
The right side is positive; normalize it to one.
Then $X^*(\psi)\leq1/\tau$.  Since
$X(v)\leq K\|v\|_1$, testing on point masses yields
$\|\psi\|_\infty\leq K/\tau$.
Also $\|\psi\|_2>1/M$, by the strict lower bound on
$\operatorname{Re}\langle b,\psi\rangle$.
Duality gives
$X(\psi)X^*(\psi)\geq\|\psi\|_2^2$ and consequently
$X(\psi)>\tau/M^2$.
The detection hypothesis contradicts
$R\sup_Q|\langle Q,\psi\rangle|\leq1$.
The indicated convex sum therefore exists with $X(e)\leq\tau$.
Its structural part has $\ell^1$ coefficient norm at most $R$.
Sample $\ell$ independent terms with probabilities proportional
to their coefficient magnitudes, absorbing their phases into
$\mathcal S$.  The expected squared $L^2$ error of the empirical
average is at most $R^2/\ell$.  Take
$\ell\geq R^2K^2/\tau^2$ and use $X\leq K\|\cdot\|_2$.
\end{proof}

In the application to shift comparison,
\begin{equation}\label{shift:testing-seminorm}
 X(e)=\E_h\sup_L
 \left|\E_ne(n)B(n+h)L(n)\right|,
\end{equation}
where $L$ ranges over complex, unit-bounded degree-$r$ niltests of
controlled complexity and all their cyclic translates.  Thus the
family includes $n\mapsto L_0(n+t)$, with addition taken in $G$,
even when its representative path wraps.  The supremum allows a
different test and translation at each $h$.  It is a seminorm and
$X(e)\leq\|B\|_\infty\|e\|_1$.
Choosing an approximately maximizing test and aligning its phase,
then applying Cauchy--Schwarz at $n+h$, gives
\begin{equation}\label{shift:detector-cs}
 X(e)^2\leq\|B\|_\infty^2
 \E_{h,k}\left|
 \E_n e(n)\overline{e(n+h-k)}
 L_h(n)\overline{L_k(n+h-k)}\right|.
\end{equation}
The product in \eqref{shift:detector-cs} contains at most two
translated representative niltests.  Lemma~\ref{prelim:dual-gowers}
therefore applies uniformly, and the derivative identity gives a
quantitative $U^{r+2}$ detector.
The quasipolynomial inverse theorem
\cite{LSSInverse2024} then supplies the structural tests for
Lemma~\ref{shift:modeling}.  These structural tests are ordinary
degree-$(r+1)$ niltests; only the testing family has been enlarged.
The translated family permits the change of variable $m=n+h$
when modeling the second input: $L_h(n)$ becomes $L_h(m-h)$.
The supremum handles these translations directly, without a union
bound over them.  All cutoff errors in the converse estimate are
chosen below the detection threshold.

\begin{lemma}[A common anchor]\label{shift:anchor}
Let $H$ be a finite nonempty set.  For $1\leq j\leq m$ and
$h\in H$, let $v_{h,j}$ be complex functions of supremum at most
one on finite probability spaces $\Omega_j$.
Suppose $|\E v_{h,j}|\geq\delta>0$ for all $h,j$.
There exist $h_0\in H$ and $H'\subseteq H$ with
$|H'|\geq\delta^{2m}|H|/2$ such that
\[
 |\E v_{h,j}\overline{v_{h_0,j}}|
 \geq\delta^{2m}/2\qquad(h\in H',\ 1\leq j\leq m).
\]
\end{lemma}
\begin{proof}
Multiply each $v_{h,j}$ by a unit scalar so its mean is real
and nonnegative, and tensor the functions in their independent
variables to form $V_h$ on $\prod_j\Omega_j$.
Then $\E V_h\geq\delta^m$ and $\|V_h\|_2\leq1$.
Consequently
\[
 \E_{h,k}\operatorname{Re}\langle V_h,V_k\rangle
 =\|\E_hV_h\|_2^2\geq\delta^{2m}.
\]
Choose $h_0$ with the same lower bound after averaging only $h$.
At least a $\delta^{2m}/2$ fraction have real inner product at
least $\delta^{2m}/2$.  For these $h$, the absolute inner product
is a product of $m$ numbers in $[0,1]$, so each factor is at least
$\delta^{2m}/2$.  Removing the aligning phases changes no absolute
value.
\end{proof}

\subsection{Recovering tests from two fixed observations}

The next lemma records precisely what positive discretization can
do after the requisite geometry has been established.

\begin{lemma}[Positive discretization]\label{shift:discretization}
Fix finite families of nonnegative functions
$\{A_i\}_{i\in I}$ and $\{B_j\}_{j\in J}$ on $G$, each summing
to one.  For each shift $h$, suppose that, outside a set $E_h$,
\[
 T_h(n)=F_h(z_h(n),\xi_h(n)),
\]
where $0\leq F_h\leq1$ and $F_h$ is $L$-Lipschitz in its second
variable.  Suppose that whenever
$A_i(n)B_j(n+h)$ and $A_i(n')B_j(n'+h)$ are positive, with
$n,n'\notin E_h$, one has
$d(\xi_h(n),\xi_h(n'))\leq\rho$.
Then there are functions
$C_{ij,h}(n)=F_h(z_h(n),\xi_{ij,h})$ such that
\begin{equation}\label{shift:positive-discretization}
 \left|T_h(n)-\sum_{i,j}A_i(n)B_j(n+h)C_{ij,h}(n)\right|
 \leq L\rho\quad(n\notin E_h).
\end{equation}
An empty cell pair may be assigned $C_{ij,h}=0$.
If the first-variable sections of $F_h$ are bounded Lipschitz
functions on a degree-$(r-1)$ nilmanifold and $z_h$ is a polynomial
path on it, then every $C_{ij,h}$ is a degree-$(r-1)$ niltest
with the corresponding uniform complexity.
\end{lemma}
\begin{proof}
For every nonempty cell pair select a point
$n_{ij,h}\notin E_h$ where its weight is positive and put
$\xi_{ij,h}=\xi_h(n_{ij,h})$.
For each contributing pair the replacement error is at most
$L\rho$.  Sum against the nonnegative weights, whose sum is one.
The final assertion follows by restricting $F_h$ to the chosen
second coordinate.
\end{proof}

\begin{lemma}[Uniform recovery after right normalization]
\label{shift:uniform-cover}
Fix controlled rational nilpotent groups, lattices, and a rational
subgroup $V$ of their product.  Fix a controlled lattice $\Lambda_V$
on $V$.  There are fixed fine covers of the input lattices with
the following property.  Fix bounded left translations and arbitrary
right translations by elements of the original input lattices, and
apply these same translations to both of two tuples.  Suppose that
both normalized tuples belong to $V$.  If the two original
tuples are within $\rho$ on the fixed input covers, then their
normalized images are within $\mathcal B(p)\rho$ on
$V/\Lambda_V$, provided $\rho\leq\mathcal B(p)^{-1}$.
The covers and constants do not depend on the right translations.
The left translations may instead differ by at most $\rho$ in each
coordinate; the same conclusion holds after enlarging its constant.
One may additionally require compatibility with right translation
by rational elements having a fixed common denominator
$q\leq\mathcal B(p)$.
\end{lemma}
\begin{proof}
First choose a fine ambient product lattice $\Lambda_0$ such that
$V\cap\Lambda_0\subseteq\Lambda_V$.
The restriction of the ambient quotient metric has a controlled
local inverse on $V/(V\cap\Lambda_0)$.
To see this directly, reduce both points using $V\cap\Lambda_0$
to bounded representatives.  An ambient lattice element witnessing
closeness is then bounded.  Its logarithm lies in a controlled
rational grid; if its distance from the rational subspace
$\log V$ is below the rational separation bound, it lies in that
subspace exactly.  It therefore belongs to $V\cap\Lambda_0$.
Controlled coordinate changes compare the ambient and intrinsic
small distances.  This is also the metric argument in
Lemma~\ref{prelim:reconstruction}.

Choose each input cover $\Lambda$ inside the corresponding factor
of $\Lambda_0$, normal in the original lattice $\Gamma$.
Here is an explicit construction.  If
$\log\Gamma\subseteq l^{-1}\Z^D$ and
$m_0\Z^D\subseteq\log\Lambda_0$, choose a controlled integer
$m$ divisible by $lm_0$, and set
\[
 \Lambda=\langle\gamma^m:\gamma\in\Gamma\rangle.
\]
Every generator lies in $\Lambda_0$ because
$\log(\gamma^m)=m\log\gamma\in m_0\Z^D$.
This subgroup is characteristic in $\Gamma$.
A Mal'cev normal form with $D$ coordinates gives
$[\Gamma:\Lambda]\leq m^D$ by reducing each exponent modulo
$m$; powers of the Mal'cev generators give controlled grids for
$\Lambda$.  Thus all logarithmic index and coordinate costs
remain polynomial.

For one coordinate, let
$a=\epsilon^{-1}x\gamma^{-1}$ and
$a'=\epsilon^{-1}x'\gamma^{-1}$, with $\epsilon$ bounded and
$\gamma\in\Gamma$.  If
$d_G(x\lambda,x')\leq\rho$ for $\lambda\in\Lambda$, then
\[
 a(\gamma\lambda\gamma^{-1})
   =\epsilon^{-1}x\lambda\gamma^{-1},\qquad
 d_G\bigl(a(\gamma\lambda\gamma^{-1}),a'\bigr)
   \leq\Lip(L_{\epsilon^{-1}})\rho.
\]
Normality puts $\gamma\lambda\gamma^{-1}$ in
$\Lambda\subseteq\Lambda_0$.  Thus normalized ambient distances
are controlled uniformly in $\gamma$.  Apply the preceding local
inverse on $V$.
If the bounded left normalizers are $\epsilon$ and $\epsilon'$,
right invariance gives
\[
 d_G(\epsilon^{-1}x'\gamma^{-1},
       (\epsilon')^{-1}x'\gamma^{-1})
 =d_G(\epsilon^{-1},(\epsilon')^{-1}).
\]
Inversion is controlled on the bounded set of normalizers, so their
variation adds at most $\mathcal B(p)\rho$ before applying the local
inverse.  No orbit representative enters this error bound.

For rational right translations, the adjoint formula has a common
denominator bounded in terms of $q$ and the structural data; its
numerators need not be bounded.  Choose the earlier congruence
strides sufficiently divisible to force the required conjugates
into $\Lambda_0$.  This requires no bound on those numerators.
Again normal refinement inside $\Gamma$ handles arbitrary lattice
right translations at the same time.
\end{proof}

\begin{proposition}[Higher-degree propagation]
\label{shift:higher}
Lemmas~
\ref{prelim:reconstruction}, \ref{prelim:symbol-splitting},
\ref{prelim:square}, \ref{prelim:dual-gowers}, and
\ref{prelim:step-drop}, together with the quasipolynomial inverse
theorem, imply the propagation principle in
Theorem~\ref{shift:required} for every fixed degree.
\end{proposition}
\begin{proof}
Induct on $d$, using Proposition~\ref{shift:degree-one} for the
base case.  Set $r=d-1\geq1$ and
$a=f-(1+\epsilon)g$.  Suppose that a set of more than $e^{-p}N$
shifts has a violating positive degree-$r$ test $T_h$.
All losses in this proof have logarithm polynomial in $p$; the
degree and the multiplicative slack are fixed.  A precision called
``sufficiently small'' below is at least $\exp(-\mathcal P(p))$.
No list is obtained by discretizing the functions $f,g,J$ or the
coefficients of a polynomial orbit.

\emph{Structural models at the two sites.}
For the seminorm \eqref{shift:testing-seminorm},
\eqref{shift:detector-cs} implies the needed detector as follows.
A lower bound $X(e)\geq\tau$, after normalizing its cap, gives
a proportion $\mathcal B(p)^{-1}$ of pairs $(h,k)$ with
derivative correlation at least $\mathcal B(p)^{-1}$.
The product of translated tests in that correlation is covered by
Lemma~\ref{prelim:dual-gowers}, including representative wraps.
Thus a proportion $\mathcal B(p)^{-1}$ of differences $h-k$
has derivative $U^{r+1}$ norm at least $\mathcal B(p)^{-1}$:
the difference is uniform, and each fixed difference has exactly
$N$ pairs.  The Gowers derivative identity gives
$\|e\|_{U^{r+2}}\geq\mathcal B(p)^{-1}$.
The quasipolynomial inverse theorem therefore gives a degree-$d$
correlator of polynomial logarithmic complexity and correlation
$\mathcal B(p)^{-1}$.

Apply Lemma~\ref{shift:modeling} with these detectors.  To fit its
compactness hypothesis, take the closure in the finite-dimensional
space of values on $G$ of the balanced family of allowed niltests.
The detector still supplies an actual niltest.  After obtaining the
finite model, approximate its finitely many closure points by
actual niltests to any prescribed further $L^2$ accuracy.
The cap bound $X\leq\|B\|_\infty\|\cdot\|_2$ controls this
additional error.  Thus a bounded function at either site has,
in the corresponding seminorm, a model with list length and
$\ell^1$ coefficient sum at most $\mathcal B(p)$.

\emph{Positive Fourier projection and fixed partners.}
Write $T_h=F_h(P_h(n)\Gamma_h)$.
Truncate its top-torus Fourier expansion at sufficiently high
controlled precision.  Declare a frequency significant if the
corresponding vertical term has correlation against $a(n)J(n+h)$
larger in absolute value than a sufficiently small controlled
threshold.  Average $F_h$ over the connected annihilator of the
rational span of the significant frequencies.
This is a positive averaging operator, so the resulting test
still lies in $[0,1]$.  It changes the tested integral by at most
the Fourier truncation error plus the number of discarded truncated frequencies
times the significance threshold.  Hence a positive violation of
size $\mathcal B(p)^{-1}$ remains.  Replace $F_h$ and $T_h$ by
these averages and retain their notation.  Each significant vertical
mode is unchanged by this replacement.

Choose a basis $\eta_{h,1},\ldots,\eta_{h,m}$ from the significant
frequencies, with corresponding vertical terms $U_{h,j}$.
After pigeonholing, $m$ is fixed and at most $\mathcal P(p)$.
If $m=0$, the averaged test factors through the degree-$(r-1)$
quotient, and degree-$r$ shift comparison already rules out the
violations after a polynomial enlargement of its budget.  Thus
assume $m\geq1$.
Model $a$ with other input $J$.  Next model $J$ for each fixed
structural term of that first model.  Choose the second precision
after the first list and its coefficient bound have been fixed.
For the latter application, write $u=n+h$ and use the shift $-h$:
the fixed structural factor is evaluated at $u-h$, and the varying
test $U_{h,j}(u-h)$ belongs to the translated testing family in
\eqref{shift:testing-seminorm}.
Markov's inequality applied to the modeling seminorm errors removes
only a negligible proportion of shifts, uniformly for every
selected $U_{h,j}$.  Expanding the two models shows that each
remaining significant term has partners $Q_j,Q'_j$ with
\begin{equation}\label{shift:partner-bias}
 \left|\E_n U_{h,j}(n)Q_j(n)Q'_j(n+h)\right|
 \geq\mathcal B(p)^{-1}.
\end{equation}
The factors can be scaled to have cap one.
Take vertical modes in the top-torus expansions of both partners,
and pigeonhole these choices for all $j$ simultaneously.
The number of choices is $\mathcal B(p)^{\mathcal P(p)}
=\mathcal B(p)$.  We retain a set of shifts of density
$\mathcal B(p)^{-1}$ on which all partner maps and functions
are fixed.

\emph{An anchor and a common rational factorization.}
Apply Lemma~\ref{shift:anchor} to the functions in
\eqref{shift:partner-bias}, tensoring in $j$.
For one fixed $h_0$, a further set of density
$\mathcal B(p)^{-1}$ satisfies, for every $j$,
\begin{equation}\label{shift:anchored-bias}
 \left|\E_n U_{h,j}(n)\overline{U_{h_0,j}(n)}
 |Q_j(n)|^2Q'_j(n+h)\overline{Q'_j(n+h_0)}\right|
 \geq\mathcal B(p)^{-1}.
\end{equation}
For each $j$, divide the representative interval at the two wrap
points in the shifted partners.  Intervals shorter than a
sufficiently small $\mathcal B(p)^{-1}N$ have negligible total
contribution.  One remaining interval has both relative length
and normalized bias at least $\mathcal B(p)^{-1}$.
On it, extend the partners using their corresponding ordinary
integer shifts.  Pigeonhole the finitely many anchor branches.

Write $Q_j(n)=V_j(g_j(n)\Gamma_j)$ on $G_j/\Gamma_j$.
Since $Q_j$ is vertical at degree $d$, $|Q_j|^2$ is a function
of the degree-$r$ quotient path $a_j(n)=g_j(n)(G_j)_d$ in
$\overline G_j=G_j/(G_j)_d$.
Let $s_{h,j}(n)$ be the normalized square path supplied by
Lemma~\ref{prelim:square} for the last two factors in
\eqref{shift:anchored-bias}, and let $S_j$ be its degree-$r$
quotient group.  Taking top Fourier modes of the functions on
$\overline G_j$ and $S_j$ selects a pair that retains bias
$\mathcal B(p)^{-1}$ on the selected interval.  Their frequencies,
denoted by $\alpha_{h,j}$ and $\beta_{h,j}$, give a unit-bounded
degree-$r$ test on
\[
 D_{h,j}(n)=(P_h(n),P_{h_0}(n),a_j(n),s_{h,j}(n))
\]
with combined top frequency
\[
 \Lambda_{h,j}
 = (\eta_{h,j},-\eta_{h_0,j},\alpha_{h,j},\beta_{h,j}).
\]
Apply Lemma~\ref{prelim:step-drop} to this biased test, obtaining
a graded rational subalgebra $W_{h,j}$ whose top layer is
annihilated by $\Lambda_{h,j}$.

Writing $G_h$ for the ambient group of $P_h$, use the common
degree-$r$ product map
\[
 (P_h,Z_h)=(P_h,P_{h_0},(a_j,s_{h,j})_{j=1}^m)
 \quad\text{into}\quad
 G_h\times G_{h_0}\times\prod_{j=1}^m(\overline G_j\times S_j).
\]
Let $\rho_j$ be its coordinate projection to the four factors
of $D_{h,j}$.  The inverse image of $W_{h,j}$ under the induced
graded projection $\operatorname{gr}(\rho_j)$ admits a splitting
of the total symbol: lift the slow and rational symbols by placing
identities in the omitted coordinates.  The selected intervals may
differ with $j$, but parameter translation leaves the symbol unchanged,
and their lengths are at least $\mathcal B(p)^{-1}N$.  Thus all
these splittings have slow bounds at scale $N$, after enlarging the
budget.  Lemma~\ref{prelim:symbol-splitting} gives a simultaneous
splitting through the intersection $U_h$ of these inverse images.
Its lift is an actual factorization
\begin{equation}\label{shift:actual-factorization}
 (P_h,Z_h)=E_hb_hR_h\kappa_h,
 \qquad b_h(n)\in H_h,
\end{equation}
with its refiltration $(H_h)_i$; $E_h$ is fully slow, $R_h$ is
rational with bounded common denominator, and $\kappa_h$ is
a constant lattice element.

Write $\pi$ for the projection to $Z_h$.  The total filtration
has degree $r$, so $(H_h)_r=\exp((U_h)_r)$, and each
$\Lambda_{h,j}\circ\mathrm d\rho_j$ annihilates $(U_h)_r$.
On
\begin{equation}\label{shift:central-kernel}
 K_h=\ker\bigl((H_h)_r\longrightarrow\pi((H_h)_r)\bigr),
\end{equation}
all coordinates except the $G_h$ coordinate are the identity.
The frequency equation therefore says that every $\eta_{h,j}$
annihilates the logarithm of the first-coordinate projection of $K_h$.
All frequencies of the projected positive function belong to their
rational span, so the function is invariant under $K_h$.
This kernel is the exponential of a rational subspace of the central
top layer, hence is connected, rational, and central in the total group.

\emph{Fixing rational data before choosing covers.}
The partner groups and their square quotients are fixed.
The controlled rational group data for the first factor, the
subgroups $H_h$, and their filtrations and quotient maps range over
at most $\mathcal B(p)$ possibilities: bases use polynomially
many rational entries, each of bounded height.
The same finite bound applies to controlled lattice choices;
one may work with the integral congruence covers constructed in
the preliminaries.  Pigeonhole all these rational choices, the
common denominators of $R_h$, and any required congruence moduli.
The retained shift density is still $\mathcal B(p)^{-1}$.
We now write $H$ for the fixed subgroup and regard all rational
diagrams as fixed.  The polynomial orbits, bounded real constants,
and Lipschitz functions may still depend on $h$.
We have not taken a common multiple over exponentially many
alternative moduli.

\emph{Freezing the slow and rational factors.}
Choose a fixed fine partition of the coordinate $n/N$ into cells
of diameter $\rho$.  On one such cell the slow factor $E_h(n)$
varies by at most $\mathcal B(p)\rho$; its values and the values
of its inverse are bounded by $\mathcal B(p)$.
By rational polynomial periodicity there is a fixed
$q\leq\mathcal B(p)$ such that, on each class $n\bmod q$,
the right lattice class needed to evaluate the first component
of $R_h(n)$ is constant.  Use a bounded representative of that
class.  After freezing these factors, the original test is,
up to $\mathcal B(p)\rho$, a controlled positive Lipschitz
function of $b_h(n)$ on a fine lattice cover of $H$.
Its invariance under \eqref{shift:central-kernel} persists,
since that subgroup is central.

Let $G_Z$ be the fixed full product group containing the other
inputs $Z_h$.  Its dimension is polynomially bounded and its
filtration has degree at most $r$: its factors are the anchor
degree-$r$ paths, the degree-$r$ quotients of the unshifted partners,
and the normalized square quotients.  Use the homomorphism
\[
 H\longrightarrow H/H_r\times G_Z,\qquad
 b\longmapsto (bH_r,\pi(b)),
\]
whose kernel is $K_h=K$, now fixed.
Lemma~\ref{prelim:reconstruction} represents the frozen function
as a positive bounded Lipschitz function of these two image
coordinates on fine product covers, at cost $\mathcal B(p)$.
Write $\Lambda\leq G_Z$ for the full ambient lattice in the second
cover.  Using $G_Z$, rather than only $\pi(H)$, is necessary because
the other factors in the projected factorization need not lie in
$\pi(H)$.  Extend from the image to the product metric space by
Lipschitz extension and clip to $[0,1]$.
The first coordinate is a polynomial path of degree $r-1$.
The remaining task is to recover the second coordinate from
fixed observations at $n$ and $n+h$.

For this recovery, also freeze the \emph{left} lattice class of
$\pi(R_h(n))$ for the full ambient target lattice $\Lambda$.
Its cocompactness supplies bounded rational representatives of
these classes.  Rational periodicity, whose period depends only on
the fixed structural data and common denominator, permits a single
refinement of $q$ by a controlled amount, uniformly over all retained
$h$ and simultaneously with the earlier right-class requirement.
On a residue class
write $\pi(R_h(n))=\lambda_n z_{h,c}$, with
$\lambda_n\in\Lambda$ and $z_{h,c}$ a bounded rational
representative.  Projecting \eqref{shift:actual-factorization}
then gives the exact identity
\begin{equation}\label{shift:left-class-formula}
 \pi(b_h(n))\Lambda
 =\pi(E_h(n))^{-1}Z_h(n)\pi(\kappa_h)^{-1}
 z_{h,c}^{-1}\Lambda.
\end{equation}
The order of these factors is essential.  A right lattice class
of $R_h$ alone would not justify this identity.

\emph{Recovering squares from fixed observations.}
The anchor paths and unshifted quotient paths in $Z_h$ are fixed
paths at $n$.  A square path is obtained from two original
partner paths by bounded left normalization and a lattice right
normalization, followed by the fixed subgroup inclusion and
square quotient.  Lemma~\ref{shift:uniform-cover} gives fixed
input covers on which this operation is uniformly Lipschitz.
In particular, the arbitrary magnitude of its lattice normalizer
does not enter the metric estimate.
The same lemma makes the right factors
$\pi(\kappa_h)^{-1}z_{h,c}^{-1}$ in
\eqref{shift:left-class-formula} compatible with fixed finer
covers.  For fixed $h$, residue class and wrap branches, both
locations being compared use the same square normalizers and the
same right factors.  The left slow inverses can differ by
$\mathcal B(p)\rho$; the varying-left-normalizer estimate in
Lemma~\ref{shift:uniform-cover} controls the resulting error
independently of the orbit representatives.

The ordinary extension of an anchored partner is a fixed path
$n\mapsto g'_j(n+h_0-b_jN)$, with its branch $b_j$ already fixed.
An ordinary extension at the shifted site is one of the fixed
paths
\[
 m\longmapsto g'_j(m+tN),\qquad t\in\{-1,0,1\},
 \quad m=\operatorname{rep}(n+h).
\]
Which branch is used is determined by $h$ and the coarse
$n/N$ cell except when that cell meets the wrap point.
Discard a neighborhood of width $O(\rho N)$ around that point
and the endpoint of the interval-coordinate chart.  This has
relative size $O(\rho+1/N)$, uniformly in $h$.
There are only polynomially many fixed paths, including every
branch required for all partners.

\emph{Positive partitions and the uniform oscillation estimate.}
At the first site measure those fixed paths on their chosen
covers, the coordinate $n/N$, and the residue $n\bmod q$.
At the second site measure all the fixed shifted-path alternatives.
Take positive Lipschitz partitions of unity into cells of
diameter at most $\rho$ on these two product spaces.
The residue partition can be realized exactly on the grid by a
circle path $n\mapsto n/q\bmod1$, with Lipschitz cost $O(q)$.
The interval-coordinate partition is realized similarly away from
the discarded endpoint neighborhood.
Since the total dimension and $\log(1/\rho)$ are polynomial in
$p$, both families have at most $\mathcal B(p)$ members and
complexity at most $\mathcal P(p)$.  Denote them by $A_i(n)$
and $B_j(n+h)$; these are fixed degree-$d$ niltests.

For a fixed $h$, consider two good locations contributing to the
same pair of cells.  Their frozen residue classes and wrap branches
agree; their slow factors differ by at most
$\mathcal B(p)\rho$.  Their fixed orbit observations are within
$O(\rho)$ on the chosen covers.
Lemma~\ref{shift:uniform-cover}, followed by
\eqref{shift:left-class-formula}, therefore gives
\begin{equation}\label{shift:uniform-oscillation}
 d\bigl(\pi(b_h(n))\Lambda,\pi(b_h(n'))\Lambda\bigr)
 \leq\mathcal B(p)\rho.
\end{equation}
This is a uniform estimate over all retained $h$, rather than
a continuity assertion whose constant may depend on $h$.

Freeze the first-site slow value separately for each cell $i$.
For every nonempty pair $(i,j)$ choose a good point and replace
the second image coordinate in the extended reconstruction
function by its value there.  The resulting function of
$b_h(n)H_r$ is a positive degree-$(r-1)$ test $C_{ij,h}$.
The proof of Lemma~\ref{shift:discretization}, allowing the
frozen function to be chosen separately for each first-site cell,
gives
\[
 T_h(n)=\sum_{i,j}A_i(n)B_j(n+h)C_{ij,h}(n)+e_h(n),
 \qquad \E|e_h|\leq\mathcal B(p)\rho+O(1/N).
\]
No sum of errors over all cell pairs is necessary, since their
nonnegative weights sum to one.
Choose $\rho$ small enough, and the length cutoff large enough,
to make this error adequate for Lemma~\ref{shift:fixed-lists-suffice}.
All costs remain within $\mathcal P(p),\mathcal B(p)$.

The retained shift set and its positive violations both have
controlled inverse-exponential lower bounds.
Lemma~\ref{shift:fixed-lists-suffice} now contradicts them by
degree-$r$ shift comparison.  Its finitely many fixed pairs are
chosen before that application; the lower tests $C_{ij,h}$ may
depend arbitrarily on $h$ because the induction is uniform in
those tests.  Choosing budgets backwards through the fixed degree
depth gives a single polynomial $P=(2+p)^{C(d,\epsilon)}$.
This proves shift comparison in every fixed degree.
\end{proof}
